\section{Problem Statement, Overview, and Objectives}

%\subsection{Problem Statement}
%\input{problem}

Differential testing has long been considered one of the most powerful
approaches to ensure the security, reliability, and correctness of
software, \emph{when applicable}.  Differential testing offers a unique
solution to the difficulty of specifying what a software system should
(and should not) do: in essence, it tests a function $f$ by checking
its behavior matches function $g$ that is \emph{an independent
  implementation of the same functionality}, and defines test failures
as cases where $f(x) \neq g(x)$.  The ``functions'' $f$ and $g$ in
realistic settings are seldom actual code-level functions, of course,
but more typically complex systems ranging from data-structure
implementations to complex parsers, cryptographic protocols~\cite{aumasson2017automated} and, especially, compilers~\cite{Differential} and file
systems~\cite{ICSEDiff}.

The ubiquity of differential approaches to testing compilers and file
systems is no accident:  these are two critical pieces of system
infrastructure where genuinely independent implementations are often
available, and standards such as ANSI C and POSIX make comparison
reasonable, if not trivial; the same applies, though less frequently,
to cryptographic algorithm and protocol implementations.  Of course, since file systems may rest on different
hardware substrates, and compilers may make differenct choices where a
language like C allows for implementation defined behavior, in practice even
in these best-case scenarios, the problem is usually not one of checking that $f(x)
= g(x)$, but of checking that $O_f(f(I_f(x))) = O_g(g(I_g(x)))$, where
$O_f$, $O_g$, $I_f$, and $I_g$ are respectively, output and input
abstraction functions for $f$ and $g$, that allow for known
differences in behavior and input spaces, and there are often also
known inputs that violate the differential specification in a known
and allowed way, making the full specification of the form $P(x)
\Rightarrow O_f(f(I_f(x))) = O_g(g(I_g(x)))$.

Despite these complexities, differential testing is extremely
attractive, especially in settings where inputs are automatically (and
stochastically) generated, such as random testing and \emph{fuzzing}.
Indeed, arguably the fundamental limitation of modern fuzzing, despite
its huge success in finding software vulnerabilities and bugs, is the
typical weakness of fuzzing oracles.  It is extremely hard to
determine if a system is behaving correctly on an arbitrary, often
extremely exotic, input.  This leads to results such as the difference
between functional tests and fuzzing harnesses for Bitcoin core, where
equivalent, almost 100\% branch and code coverage yields huge
discrepancies in mutation detection, with functional tests detecting
almost 80\% of injected bugs while fuzzing, despite perhaps more
aggressive exploration of the state space, detects only 12\% of these
bugs~\cite{seip2022}.  A differential specification makes it possible for fuzzing to
detect subtle semantic errors, not just memory safety violations and
assertion failures.

It seems likely that the importance of differential fuzzing will only
grow in the future, for two reasons, both deriving from the use of
LLMs and agentic coding.  First, the primary reason differential
fuzzing is seldom used in practice is that the cost of an independent
implementation of a software system, even in security or safety
critical settings, is usually prohibitive.  This cost is especially high given
that previous work has shown that ``independent'' implementations are
seldom as independent (in terms of correlation of bugs) as desirable~\cite{KNIGHT1985159}.
These means that differential testing via separate implementations of
functionality is usually only performed  in cases of extremely complex specifications implemented by
radically separated teams (e.g., compilers and file systems) with
orthogonal incentives to produce their own implementations, or for very
simple specifications such as data structures.  Code synthesis by AI
models sidesteps this problem:  once a specification that produces a
single workable implementation is available, having a single model
produce multiple implementations is often relatively cheap.  Using
multiple models or changing prompts to focus code generation on, for
example, a faster implementation, an implementation with a smaller
memory footprint, a more readable implementation, an implementation
using an unusual algorithm, and so forth, seems likely to produce
controlled divergence that may exceed that of human developers.  If
the development cost of multiple implementations is radically reduced,
the advantages of differential specification are clear.

One pre-AI way to work around the cost of a second implementation, of
course, is to use differential fuzzing in a \emph{regression} setting,
comparing a new version of software to an older, better-tested
version.  However, in practice, new versions of a software system are often introduced precisely
because the functional interface has changed in some way, and the work
required to conceive, implement, and maintain a mapping between new
and old versions (which will be less stable, as each new version tends to
remove some parts of the current mapping, as the ``old'' version is
now the version that was new, as well as add new changes) discourages
use in this setting, especially as the bug correlation problem reduces
the power of differential testing.  Again, LLMs are probably
well-suited to constructing and maintaining version mappings at low
cost\footnote{We believe this wil be especially true if LLMs have
  access to a high-quality DSL for constructing differential
  harnesses, which is one of the goals of this proposal.}.

Second, differential approaches seem ideal for probing ambiguities in
specifications presented to AIs, and AI misunderstandings:  if
multiple implementations are produced from a specification, and
checked for equivalence, supervisors of AI-driven development can
examine cases where the AI's ``view'' of the software's behavior is
unstable, and can work to remove ambiguity and correct mistakes.

Differential fuzzing is, of course, already an important practice,
particularly in security-critical code development.  However,
differential fuzzing as practiced now is often limited in
effectiveness.  All too often, the approach taken is to independently
fuzz two implementations, then compare their behavior over the union
of the produced corpuses of interesting inputs.  While this can
discover bugs, it is effectively throwing away most of the power of
fuzzing, in two important ways.  First, the implementations are only
tested differentially for a small number of inputs, the final results
of fuzzing runs on the two separate systems; this missed out on the
opportunity to detect divergence and bugs on potentially millions or
billions of executions.  Second, the two independent corpuses are only
derived from efforts to cover behaviors of the two systems
independently; if an input does not produce novelty in terms of either system
in isolation, but represents a unique \emph{pair of coverage patterns} for a
shared input, this method is simply incapable of identifying and
exploiting such inputs.  Such inputs, however, are, intuitively, the
most promising in terms of detecting divergences in implementation behaviors.

With somewhat more effort, users of differential fuzzing can compile two
implementations in such a way that a fuzzer can see the combined
coverage of both systems, essentially a program path over the merged program{\tt r1 =
  f(x); r2 = g(x); assert (r1==r2);}.  This overcomes the above
objections, but still in fact leaves much to be desired, in terms of
taking advantage of the differential nature of the specification.

\subsection{First-Class Differential Fuzzing}

In computer science, fully leveraging an approach tends to require
making it truly \emph{first-class}, supported by language and analysis
tools; without OO constructs or native-language threads, objects and
concurrency are limited in applicability, and the recent jump in
sophisticated usage of functional programming techniques in mainstream
languages can essentially be traced to the widespread availability of
real lambdas and closures.

Similarly, differential specification (arguably the most powerful and
``usable'' of all specification techniques) can only be fully
exploited when \emph{fuzzing algorithms} and \emph{specification
  techniques} are differential-aware.  This requires abstracting away
from the specifics of a particular instance of differential
testing/fuzzing.  The most important previous work along this line,
the NEZHA framework of Pesios et al.~\cite{nezha} offered a start on
such an approach, and proved extremely successful (finding differences
between implementations at a rate 20-50 times higher than previous
approaches) in practice, but failed to truly offer up first-class
differential fuzzing.

In order to understand the limits of current approaches, it is useful
to think about a very simplified description of what a modern
mutation-based fuzzer such as AFL, libFuzzer, or Centipede, does, at a
high level that essentially applies to all fuzzers:

\begin{enumerate}
\item First, a fuzzer selects an input, a \emph{seed} that it has already executed
 (possibly presented to it as part of an initial \emph{seed corpus} of
  interesting inputs, e.g., valid media files when fuzzing a media
  parser), based on a \emph{seed scheduling} algorithm that selects
  the most promising such inputs to work with.  At this stage the
  fuzzer may also compute an \emph{energy}: how much time to spend
  working on this particular seed (energy can also be modified once
  the fuzzer starts seeing results of working with the seed, but we
  are offering a simplified version of the process here).
\item Then, repeatedly, up until the limits set by \emph{energy}, the
  fuzzer \emph{mutates} the seed, making a randomly chosen change to
  the input; these changes can be simple byte-based mutations (often
  referred to as ``havoc'' operators) or more targeted and
  domain-specific, as in some compiler-fuzzing tools~\cite{cc2022}.
\item For each such modified variant of the seed, the fuzzer executes
  the mutant of the seed; if a crash is detected (in differential
  fuzzing, the executable should be set up so that divergence in
  behaviors triggers a crash) the input is preserved as representing a
  potential bug.  If the input does not cause a crash, the fuzzer
  determines if the input triggers \emph{novel} behavior.  The
  definition of novelty can vary; some fuzzers only consider inputs
  that cover new branches or basic code blocks novel (this is
  libFuzzer's default mode, for example); other fuzzers (AFL and
  family, for example) make use of a more complex edge-hit-count
  pattern approach that can approximate a form of path coverage.  If
  the input is novel, it is added to the set of seeds to schedule and
  spend energy on collected by the fuzzer, and the loop eventually
  returns to step 1, when energy for the current seed is exhausted.
  The identification of novel inputs is what lets the fuzzer slowly
  accumulate new interesting program behaviors and deeply explore the
  program, searching for bugs.
\end{enumerate}

First, NEZHA's fundamental contribution was making a tuple
representation of coverage across compared methods,
$\langle \text{cov}(f), \text{cov}(g) \rangle$. This is a huge advance on
simply attempting to cover both implementations independently, then
compare over interesting paths for either approach; it captures the
importance of asymmetry and unusual \emph{pairings} of coverage paths
between compared implementations.  However, this is in essence already
done by AFL/AFL++ and other fuzzers, if the differential targets are
compiled into a single instrumented binary, and NEZHA's experimental
evaluation unfortunately only compares to
application of AFL to generate inputs based on a single
implementation.  NEZHA's $\delta$-based measure of input novelty is essentially based on a flattening
of output (in black-box settings) or program path (in grey-box
settings) tuples, so that novel pairings of $f$ and $g$ behaviors are
observable; when a single executable with a global path is produced
from the differential fuzzing targets in AFL/AFL++, the same result is
essentially achieved.  The advantage NEZHA shows over an
implementation of what they call ``global coverage'' in  libFuzzer is due
to libFuzzer's simpler novel branch/block heuristic for determining
which inputs are interesting.

However, despite missing the fact that their grey-box approach to
$\delta$-based novelty is likely simply equivalent to a
single-compiled-binary AFL/AFL++ approach, the NEZHA paper shows how
to \emph{take differential fuzzing seriously}.  This proposal aims to
go beyond this approach by modifying fuzzer internals (implemented on
top of the industry-standard AFL++ platform) to take advantage of the
differential problem in all stages of the basic fuzzing algorithm
presented above:

\begin{itemize}
  \item First, and most important, we aim to extend the notion of
    novelty to preserve, and thus explore from, inputs that do not find novel
    behavior of either $f$ or $g$ alone, but which (going beyond NEZHA and
    AFL-family methods) expose cases where either of $f$ or $g$
    \emph{changes behavior} more in response to a shared change in an input
    seed.  Intuitively, cases where one implementation's behavior is
    much more affected by a change to an input than the other are ``on
    the path'' to exposing full behavioral divergence.  This requires
    a new approach to novelty evaluation not previously explored in
    the fuzzing literature to our knowledge (relative/response-based
    novelty), where the behavior of a mutated seed is considered
    \emph{relative to the original seed behavior}.
    \item Second, the extent of such changes can be used to inform
      scheduling and energy decisions; the more input changes
      differently affect behavior, the more the input merits urgent
      investigation and substantial computation resources.
     \item Third, these basic ideas can be further extended in the
       case where $f$ and $g$ share a common code base, since $f$ is a new
       version of $g$, not an independent implementation.  In these
       cases program edges can be mapped into sets of shared and
       differing code locations, and the structure of the code changes
       can be exploited to better focus on the changes made to $g$.
       \item Finally, the fuzzer can identify when certain byte
         locations (or higher-order structure entities, in theory) or
         mutation operators are especially prone to produce large
         divergences in response for $f$ and $g$ and spend more time
         applying those more-productive mutants.
         \end{itemize}

\section{Intellectual Merit}
The essence of this proposal is the hypothesis that software security
and reliability could be greatly enhanced by development of concepts,
algorithms, and tools for exploiting the similarity of many software
systems.  The case of comparison of different versions of the same
system is the most obvious of these, where differential testing is a
kind of abstraction of regression testing; the most important case for
software security is the large body of fundamental systems software
that implements functionality that is repeatedly implemented because
it is so foundational to a variety of software systems.  Such
``ubiquitous'' software artifacts (compilers, file systems,
cryptographic algorithms, regular expression tools, data structures)
are notably also frequently involved in particularly dangerous
software vulnerabilities, because they are so widely used in different
systems, and invoked in low-level code including core security
functionality and operating system settings that operates with
elevated privileges or direct impact on security properties of a system.

\subsection{First-Class Differential Fuzzing}


In order to understand the limits of current approaches, it is useful
to think about a very simplified description of what a modern
mutation-based fuzzer such as AFL++~\cite{AFLplusplus-Woot20},
libFuzzer~\cite{libFuzzer}, or Google's Centipede~\cite{centipede}, does, at a
high level that essentially applies to all fuzzers:

\begin{enumerate}
\item First, a fuzzer selects an input, a \emph{seed} that it has already executed
 (possibly presented to it as part of an initial \emph{seed corpus} of
  interesting inputs, e.g., valid media files when fuzzing a media
  parser), based on a \emph{seed scheduling} algorithm that selects
  the most promising such inputs to work with.  At this stage the
  fuzzer may also compute an \emph{energy}: how much time to spend
  working on this particular seed (energy can also be modified once
  the fuzzer starts seeing results of working with the seed, but we
  are offering a simplified version of the process here).
\item Then, repeatedly, up until the limits set by \emph{energy}, the
  fuzzer \emph{mutates} the seed, making a randomly chosen change to
  the input; these changes can be simple byte-based mutations (often
  referred to as ``havoc'' operators) or more targeted and
  domain-specific, as in some compiler-fuzzing tools~\cite{cc2022}.
\item For each such modified variant of the seed, the fuzzer executes
  the mutant of the seed; if a crash is detected (in differential
  fuzzing, the executable should be set up so that divergence in
  behaviors triggers a crash) the input is preserved as representing a
  potential bug.  If the input does not cause a crash, the fuzzer
  determines if the input triggers \emph{novel} behavior.  The
  definition of input novelty can vary; some fuzzers only consider inputs
  that cover new branches or basic code blocks novel (this is
  libFuzzer's default mode, for example); other fuzzers (AFL and
  family, for example) make use of a more complex edge-hit-count
  pattern approach that can approximate a form of (acyclic) path coverage~\cite{ISSTA13}.  If
  the input is novel, it is added to the set of seeds to schedule and
  spend energy on collected by the fuzzer, and the loop eventually
  returns to step 1, when energy for the current seed is exhausted.
  The identification of novel inputs is what lets the fuzzer slowly
  accumulate new interesting program behaviors and deeply explore the
  program, searching for bugs.
\end{enumerate}

First, NEZHA's fundamental contribution was making a tuple
representation of coverage across compared methods,
$\langle \text{cov}(f), \text{cov}(g) \rangle$. This is a huge advance on
simply attempting to cover both implementations independently, then
compare over interesting paths for either approach; it captures the
importance of asymmetry and unusual \emph{pairings} of coverage paths
between compared implementations.  However, this is in essence already
done by AFL/AFL++ and other fuzzers, if the differential targets are
compiled into a single instrumented binary, and NEZHA's experimental
evaluation unfortunately only compares to
application of AFL to generate inputs based on a single
implementation.  NEZHA's $\delta$-based measure of input novelty is essentially based on a flattening
of output (in black-box settings) or program path (in grey-box
settings) tuples, so that novel pairings of $f$ and $g$ behaviors are
observable; when a single executable with a global path is produced
from the differential fuzzing targets in AFL/AFL++, the same result is
essentially achieved.  The advantage NEZHA shows over an
implementation of what they call ``global coverage'' in  libFuzzer is due
to libFuzzer's simpler novel branch/block heuristic for determining
which inputs are interesting.

However, despite missing the fact that their grey-box approach to
$\delta$-based novelty is likely simply equivalent to a
single-compiled-binary AFL/AFL++ approach, the NEZHA paper shows how
to \emph{take differential fuzzing seriously}.  This proposal aims to
go beyond this approach by modifying fuzzer internals (implemented on
top of the industry-standard AFL++ platform) to take advantage of the
differential problem in all stages of the basic fuzzing algorithm
presented above:

\begin{itemize}
  \item First, and most important, we aim to extend the notion of
    novelty to preserve, and thus explore from, inputs that do not find novel
    behavior of either $f$ or $g$ alone, but which (going beyond NEZHA and
    AFL-family methods) expose cases where either of $f$ or $g$
    \emph{changes behavior} more in response to a shared change in an input
    seed.  Intuitively, cases where one implementation's behavior is
    much more affected by a change to an input than the other are ``on
    the path'' to exposing full behavioral divergence.  This requires
    a new approach to novelty evaluation not previously explored in
    the fuzzing literature to our knowledge (relative/response-based
    novelty), where the behavior of a mutated seed is considered
    \emph{relative to the original seed behavior}.
    \item Second, the extent of such changes can be used to inform
      scheduling and energy decisions; the more input changes
      differently affect behavior, the more the input merits urgent
      investigation and substantial computation resources.
     \item Third, these basic ideas can be further extended in the
       case where $f$ and $g$ share a common code base, since $f$ is a new
       version of $g$, not an independent implementation.  In these
       cases program edges can be mapped into sets of shared and
       differing code locations, and the structure of the code changes
       can be exploited to better focus on the changes made to $g$.
       \item Finally, the fuzzer can identify when certain byte
         locations (or higher-order structure entities, in theory) or
         mutation operators are especially prone to produce large
         divergences in response for $f$ and $g$ and spend more time
         applying those more-productive mutants.
       \end{itemize}

 Fundamentally, these ideas can be summarized as taking the fact that
 differential fuzzing differs from most fuzzing in that it has a
 mapping between two pairs of behaviors to exploit seriously.  In
 order to make this concept clear, it is easiest to focus on the first
 proposed algorithmic change, which seems most novel and most
 promising.

 Consider two functions $f$ and $g$ that implement a critical cryptographic function,
 where the $g$ is a more complex, but more efficient,
 implementation of the same algorithm.  The specification here is simple:  the two
 functions should, given the same inputs (e.g., keys and ciphertext,
 or protocol message), yield
 the same outputs (plaintext, protocol configuration structure).
 Now, consider the behavior of a traditional AFL-style harness calling
 both functions on a set of inputs A, B, C, D, and E.  Suppose the
 harness is constructed such that first $f$ is called, then $g$, and
 the paths taken by the functions, where $f_1$ etc. are path edges traversed
 in the respective functions are:

 \begin{tabular}{l|r|r}
   Input & $f$ path & $g$ path\\
   \hline
   A & $f_1, f_4, f_6$ & $g_1, g_3, g_4, g_5$\\
       B & $f_1, f_4, f_6$ & $g_1, g_2, g_4, g_5$\\
           C & $f_1, f_2, f_3$ & $g_1, g_3, g_4, g_5$\\
               D & $f_1, f_2, f_3$ & $g_1, g_3, g_4, g_6$\\
                   E & $f_1, f_4, f_6$ & $g_3, g_2, g_8, g_9$\\
 \end{tabular}

 All of these inputs will be saved as novel, since they all execute an
 overall new sequence of edges.  This is the case despite the fact
 that input C, for example, does not produce novel behavior in either
 $f$ or $g$ alone, it is only novel with respect to both functions.
 That such sequences need to be seen as novel for effective
 differential fuzzing is a key insight of NEZHA.

 Now consider the input F.  F produces the path $f_1, f_4, f_6$ $g_1, g_3, g_4, g_5$.  Should F be saved as an important input to explore
 from?  AFL++ or the NEZHA approach would say no; F produces the same
 path through $f$ and $g$ as A, the very first input explored.
 However, this ignores how F was derived:  F is the result of mutating
 input E.  The small change made to the input did not change the
 behavior of function $f$ but dramatically changed the behavior of
 function $g$.  Had F been produced from D, on the other hand, the
 beahvior of \emph{both} $f$ and $g$ would be seen to be changed, and the
 previously seen path therefore uninteresting.  This potentially highly significant fact is invisible
 unless the fuzzer is aware of the structure of the fuzzing problem,
 and both distinguishes $f$ and $g$ behaviors \emph{and} makes use of
 the \emph{provenance} of fuzzing inputs, so that behavior change or lack
 of change can be \emph{observed} and \emph{exploited} as a signal of
 behavioral divergence.  

 All of the above proposed differential adaptations of the core
 mutation-driven fuzzing algorithm are derived from this insight.  For
 instance, the insight behind the third method, which tries to exploit
 cases where $f$ and $g$ share a largely identical code base, is that
 when a change to an input produces changes in the \emph{shared} code
 that does not differ between $f$ and $g$ this is likely to be highly
 significant:  behavior has flowed out of the non-shared portions of
 the code base to influence the behavior in ostensibly identical
 parts of the implementation.  Changes wholly contained in divergent
 parts of the code may be due to accidental aspects of the way the
 code behaves; this seems,
 intuitively, less likely when there is data flow of difference to
 shared code.  If we simplify the model by assuming the changed code
 is contained in a single block, modifying a unit of functionality
 that is internally modified but intended to have the same black box
 behavior, we have identified a case where the two black boxes must be
 behaving differently, since code outside the box has changed
 behavior.  The conclusion is less confident with more complex
 relationships between changed versions, but the basic argument for
 elevating the signal on divergence in shared code seems at minimum
 deserving of empirical exploration.

 \subsection{First-Class Differential Specification}

Lack of fuzzer awareness of differential structure contributing to
weakness in finding differentially-defined bugs is almost certainly
\emph{not} the primary reason differential testing is not more widely used, even ignoring the cost of developing independent
implementations (which may now be mitigated by AI code synthesis).
Many developers do not use differential testing methods
because in practice, the problem is seldom one of comparing two
implementations with perfectly identical behavior, and the effort required to
make differential testing actually useful can seem overwhelming.  While
the first prong of this proposal will, we believe, greatly improve the
effectiveness of differential testing for those already using it, it
is unlikely to be widely used without accompanying advances in \emph{differential specification}.

Gulzar et al. in 2019 reported on the (large and growing, at that
time) adoption of differential testing at Google~\cite{PercepDiff}.
Their results, based on a survey of 117 software engineers, in-depth
interviews with 5 experienced users of differential testing, and
analysis of over 2,000 differential tests, showed two key problems
with differential testing.  First, differential tests, as practiced,
usually based on data drawn from end-to-end execution of production
systems, were noted as slow.  This problem is in part addressable by better
differential fuzzing, which can produce fast-executing tests that
concentrate on difference-exposing tests.  The second most common
problem was that manual comparison of results is tedious and
difficult, and clear pass/fail results are hard to obtain.  This
problem is therefore significant in the setting of the study, manual difference
testing between different versions of the same software system.  It is
much more important in a differential fuzzing context, where hundreds
or thousands of false positives may be generated, behavior is more
likely to be found in obscure space of the system state space, and the
comparison may be between truly independent implementations of a standard
functionality, rather than versions of the same software system.

Differential testing harnesses such as that for the widely used Pyfakefs Python
mock file system~\cite{pyfakefs,TSTLHarness} often devote up to half their
code to allowance for divergences between a tested and reference
system; the POSIX standard is, like many other important standards,
partly nondeterministic.  The intellectual effort required to construct such
abstraction code and determine which divergences are important and
which are spurious is a major burden on would-be users of differential
testing, even when analysis reveals the categorization to be, in some
sense, ``obvious.''

Users of differential fuzzing can be divided, roughly, into two
classes:

\begin{itemize}
  \item Users expecting a high degree of automation, who do not want
    to invest a large human effort in constructing a differential specification.  In such cases, often but not always focusing on
    different versions of the same software system, users are focused
    on ease-of-use of differential testing, not on its raw power.  These users may be most likely to be developers of
    the code under test, and aim to use differential testing in large
    part to detect impact of code changes they have made outside the
    scope of their own unit tests and module~\cite{PercepDiff}.
    \item Users who are willing to devote considerable effort to
      producing a differential specification.  These cases tend to be
      focused on using differential testing as a means for detecting
      very subtle bugs and obscure behaviors, not using differential
      testing primarily as an additional, easy-to-apply, extension of
      unit tests.  File system and cryptographic algorithm
      differential testing and other high-stakes systems software
      where non-developer testing experts may be involved is a
      paradigmatic instance of this kind of high-effort differential
      testing.
    \end{itemize}

These two modes of differential testing/fuzzing naturally suggest two
modes of specification.  In the first case, the ideal implementation
of differential fuzzing takes as input two binaries, and compares
their behavior, with no additional work required by the user.  In some
cases, this is sufficient.  More commonly, however, some complexity
remains:  two systems may produce output on a channel such as a file
or stdout/stderr that is almost identical, but diverges in, e.g., the
content of text messages.  In other cases, behavior may not be so
simply observable, and users may want to compare less visible aspects
of behavior, such as system call sequences, maximum memory footprint,
or execution time.  Of course, an AFL++ extension can be developed
that allows users to specify such behavioral comparison standards on
the command line, including potentially inline regexps or
configuration files to control allowed divergence in text-based
outputs, but the triage problem of understanding divergences is a
major burden on users.  We propose to not only implement and classify
an extensive set of composable, customizable behavioral comparison
functions, but develop automated calibration methods that use the
ability of fuzzing to produce a statistically large number of
behaviors to estimate a good comparison function.  Inferring a change
in output message from a large body of stdout traces that differ only
by a regexp substitution is the kind of task that can and should be
automated, and a configuration generator can reduce the risk of such
an approach by presenting a summary of consistent differences to
users in an easily-understood format, e.g.:

\begin{code}
- Version B consistently replaces``FILE OR DIRECTORY'' with 
``FILE/DIRECTORY'' in stderr messages
- Version B consistently reorders the {\emph rename} and {\emph close}
POSIX calls from {\emph rename --> close} to {\emph close --> rename}
- Version B omitted a call to \emph{clock\_gettime()} from 95\% of
all observed executions; a set of representatives of the 5\% of cases
still calling this function were saved to the {\tt exceptions/} directory
\end{code}

The second case is where the user of differential fuzzing aims to
put in the effort to construct a harness that controls inputs to avoid
areas of known divergence, to introduce complex behavioral pairings
(e.g., where fault injection is performed in a simulated flash file
system, and this has to be handled by intercepting and modifying calls
to a ``real'' file system appropriately~\cite{ICSEDiff}), or where the
mapping between inputs and outputs to the two systems is far too
complex to infer from statistical observation, and must be modeled by
a domain expert.

WRITE MORE ABOUT THIS PART, BRING IN TSTL

Conceptually, the key feature of differential specification is that
the core of the problem is \emph{avoiding false positives}: in
traditional, assertion-based testing, or even property-based testing,
specification effort is usually largely invested in defining failure; in differential
testing, the effort is in avoiding reporting spurious failures, since
the default equality relation defines all but a very small set of
executions as failing, with no additional effort.

\subsection{PI Qualifications}

PI Groce has a long history with software analysis and verification,
including involvement in design and development of well-known fuzzing
and testing tools, including the TSTL testing tool for Python, and the DeepState
unit-fuzzing interface for C++ and C code; he is at present lead developer and
maintainer for multiple open source testing and fuzzing tools, with over 100
GitHub stars (and, in one case, more than 800 stars); he has, in the
past, managed releases of industrial open source testing tools with
more than 3000 GitHub stars.  
Groce typically reports 50-100
bugs detected using such tools to open source projects each year, and
has been awarded significant bug bounties for security-critical
compiler bugs detected using his work.  Groce was twice employed by
the Bitcoin Core team to investigate the effectiveness of their
fuzzing framework, a key component in ensuring the security of the
leading implementation of the most important cryptocurrency; this
effort included evaluations of the Bitcoin Core fuzzing team's own
internal work on differential fuzzing of Bitcoin versions.  His work
has consistently involved a marriage of theory of testing and
analysis with practical work oriented towards real-world
developers, implemented in working tools (all, as will be the case with this
project, released as open source).  Groce is author of one of the
most-cited papers on differential testing~\cite{ICSEDiff} (the 7th
most cited paper on the topic, according to Google scholar) and the TSTL testing
DSL for Python was the first (and, to our knowledge, thus far
\emph{only}) such language to include automation of
differential harness construction, including complex mappings on inputs and
outputs, as a first-class construct in declarative testing~\cite{ISSTA15,tstl}.
